% %%%%%%%%%%%%%%%%%%%%%%%%%%%%%%%%%%%%%%%%%%%%%%%%%%%%%%%%%%%%%%%%%%%%%%%%%%%%%%%%%
%
%	cap#.tex 
%  -----------
%
%	Neste arquivo você escreve o capítulo. Deve sempre iniciar com o ambiente
%	\chapter{Nome do Capitulo}
%
% %%%%%%%%%%%%%%%%%%%%%%%%%%%%%%%%%%%%%%%%%%%%%%%%%%%%%%%%%%%%%%%%%%%%%%%%%%%%%%%%%
%
\chapter{Diretrizes da pesquisa}

...

% ---
\section{Tema}
% ---

Deformações induzidas pelo processo construtivo defasado de túneis gêmeos profundos com galerias transversais considerando o comportamento instantâneo e diferido no tempo.

% ---
\section{Problema de pesquisa}
% ---

?

% ---
\section{Objetivos}
% ---

O objetivo principal do presente trabalho é...

Além disso, os objetivos secundários são:

\begin{alineas}
	
	\item objetivo secundário 1; 

	\item objetivo secundário 2;
	
	\item objetivo secundário 3.
	
\end{alineas}

% ---
\section{Delimitações}
% ---

Adota-se as seguintes delimitações para este trabalho:

\begin{alineas}
	
	\item Os túneis em estudo são considerados profundos, ou seja, são desprezadas as cargas provenientes de estruturas da superfície e também recalques induzidos pela escavação de túneis.
	
	\item O maciço pode apresentar descontinuidades e possuir um comportamento muito complexo. Portanto, para reduzir essa complexidade e mantendo a efetividade das simulações, neste trabalho admite-se um meio contínuo, homogêneo e isotrópico.
	
	\item Considerando a complexidade do estado de tensões internas do maciço devido à anisotropia, heterogeneidade e descontinuidade dos materiais, adotou-se um estado de tensões geostático-hidrostático, onde as tensões verticais e horizontais são iguais.
	
	\item velocidade de escavação constante
	
	\item modelo único de revestimento e com espessura constante 
	
	\item escavação é de face inteira
	
	\item citar os modelos utilizados (elastico, elastoplastico, viscoplastico...)
	
	\item evolução das deformações quase-estática (não se considera excitações dinâmicas)
	
\end{alineas}

% ---
\section{Metodologia/Delineamento e etapas do estudo}
% ---

Como seguirá essa pesquisa 





